\subsection{Dominant orthogonal sequences}
In \cite{jep93Kos12}, Kostant introduced his so-called ``chain cascade'' of orthogonal roots. Here we will need a version of his algorithm where we impose that all the roots of the chain cascade have a positive coefficient on $\zeta$. Note that a similar algorithm is used in \cite{jep93BM15}.

We define an integer $q$ and, for any integer $i$ such that $1\leqslant i\leqslant q$, a root $\alpha_i$ together with the subset $\Pi_i\subset\Pi$, by the following inductive process:
\begin{itemize}
\item We let $\Pi_1=\Pi$
\item Assuming that $\alpha_1,\ldots,\alpha_{i-1}$ and $\Pi_1,\ldots,\Pi_i$ have been defined, we let $\alpha_i$ be the highest root of the root system $R(\Pi_i)$ generated by $\Pi_i$
\item We let $\Sigma_i\subset\Pi_i$ be the set of simple roots $\beta$ such that $\langle\alpha_i^\jepPubCoroot,\beta\rangle\neq0$
\item If $\zeta\in\Sigma_i$, then $q=i$ and the algorithm terminates. Otherwise, $\Pi_{i+1}$ is the connected component of $\Pi_i\smallsetminus\Sigma_i$ containing $\zeta$
\end{itemize}
\begin{definition}
If $(\alpha_i)_{1\leqslant i\leqslant q}$ is the sequence defined by this process, we say that it is the \emph{maximal dominant orthogonal sequence} for $\varpi$. More generally, the sequences $(\alpha_i)_{1\leqslant i\leqslant r}$ for $1\leqslant r\leqslant q$ are called the \emph{dominant orthogonal sequences}.
\end{definition}
\begin{example}
In type $A_{n-1}$ with the standard base $(\beta_1,\ldots,\beta_{n-1})$, and for $\zeta=\beta_p$ with $2p\leqslant n$, we get $\Pi_i=\{\beta_i,\ldots,\beta_{n-i}\}$ and the maximal dominant orthogonal sequence is $(\sum_{i=1}^{n-1}\beta_i,\sum_{i=2}^{n-2}\beta_i,\ldots,\sum_{i=p}^{n-p}\beta_i),$ so $q=p$.
\end{example}
For later use, we record in Table~1 below what are the dominant orthogonal sequences $(\alpha_1,\ldots,\alpha_r)$ in all cases. The number $r$ varies from $1$ to $r_M$, see Proposition~2.19~(2).

The following proposition essentially adapts the results of \cite{jep93Kos12} to our context and explains our terminology.
\begin{proposition}
Let $(\alpha_i)_{1\leqslant i\leqslant q}$ be the maximal dominant orthogonal sequence for $\varpi$. Then
\begin{enumerate}[label=(\arabic*)]
\item The roots $\alpha_i$ are long and strongly orthogonal
\item $q=r_M$
\item For any integer $i\leqslant r_M$, $\alpha_1^\jepPubCoroot+\cdots+\alpha_i^\jepPubCoroot$ is a dominant coweight: for all $\beta$ in $\Pi$, we have $\langle\beta,\alpha_1^\jepPubCoroot+\cdots+\alpha_i^\jepPubCoroot\rangle\geqslant0$
\item $\varpi-\alpha_1-\cdots-\alpha_{r_M}$ is the lowest weight $\mu_0$ of the irreducible $G$-module $\jepPubBbb{E}$
\end{enumerate}
\end{proposition}

\begin{table}[htbp]\centering\footnotesize
\setlength{\tabcolsep}{2pt}\renewcommand{\arraystretch}{1.3}
\begin{tabular}{|c@{\quad}c|c|c|}\hline
$(G,\varpi)$ & $\mathfrak g_{\mathbb R}$, $r_M$ & $\alpha_1^{\jepPubCoroot}+\cdots+\alpha_r^{\jepPubCoroot}$ & $\alpha_r$\\\hline
\multirow{2}{*}{\begin{tabular}{c}$(A_{n-1},\varpi_p)$, $2p\leqslant n$\\[2pt]\jepGraph{0/0/0/{}/{1}/0,1/1.55/0/{}/{}/0,2/2.55/0/{}/{p}/1,3/3.55/0/{}/{}/0,4/5.1/0/{}/{n-1}/0}{1/2,2/3}{0/1,3/4}{}{}\end{tabular}} & \multirow{2}{*}{\begin{tabular}{c}$\mathfrak{su}(p,n-p)$\\[4pt]$r_M=p$\end{tabular}} & \begin{tabular}{r@{\;}c}$\text{if }2r<n:$ & \jepGraph{0/0/0/{0}/{}/0,1/1.55/0/{0}/{}/0,2/2.55/0/{1}/{r}/0,3/3.55/0/{0}/{}/0,4/5.1/0/{0}/{}/0,5/6.1/0/{1}/{n-r}/0,6/7.1/0/{0}/{}/0,7/8.65/0/{0}/{}/0}{1/2,2/3,4/5,5/6}{0/1,3/4,6/7}{}{}\end{tabular} & \jepGraph{0/0/0/{0}/{}/0,1/1.55/0/{0}/{}/0,2/2.55/0/{1}/{r}/0,3/4.1/0/{1}/{n-r}/0,4/5.1/0/{0}/{}/0,5/6.65/0/{0}/{}/0}{1/2,3/4}{0/1,2/3,4/5}{}{}\rule[-12pt]{0pt}{36pt}\\
\cline{3-4}
 & & \begin{tabular}{r@{\;}c}$\text{if }2r=n:$ & \jepGraph{0/0/0/{0}/{}/0,1/1.55/0/{0}/{}/0,2/2.55/0/{2}/{p}/0,3/3.55/0/{0}/{}/0,4/5.1/0/{0}/{}/0}{1/2,2/3}{0/1,3/4}{}{}\end{tabular} & \jepGraph{0/0/0/{0}/{}/0,1/1.55/0/{0}/{}/0,2/2.55/0/{1}/{p}/0,3/3.55/0/{0}/{}/0,4/5.1/0/{0}/{}/0}{1/2,2/3}{0/1,3/4}{}{}\rule[-12pt]{0pt}{36pt}\\
\hline
\multirow{2}{*}{\begin{tabular}{c}$(B_n,\varpi_1)$\\[2pt]\jepGraph{0/0/0/{}/{1}/1,1/1/0/{}/{}/0,2/2.55/0/{}/{}/0,3/3.55/0/{}/{}/0,4/4.55/0/{}/{n}/0}{0/1,2/3}{1/2}{3/4/1}{}\end{tabular}} & \multirow{2}{*}{\begin{tabular}{c}$\mathfrak{so}(2,2n-1)$\\[4pt]$r_M=2$\end{tabular}} & \begin{tabular}{r@{\;}c}$\text{if }r=1:$ & \jepGraph{0/0/0/{0}/{}/0,1/1/0/{1}/{}/0,2/2/0/{0}/{}/0,3/3.55/0/{0}/{}/0,4/4.55/0/{0}/{}/0,5/5.55/0/{0}/{}/0}{0/1,1/2,3/4}{2/3}{4/5/1}{}\end{tabular} & \jepGraph{0/0/0/{1}/{}/0,1/1/0/{2}/{}/0,2/2.55/0/{2}/{}/0,3/3.55/0/{2}/{}/0,4/4.55/0/{2}/{}/0}{0/1,2/3}{1/2}{3/4/1}{}\rule[-12pt]{0pt}{36pt}\\
\cline{3-4}
 & & \begin{tabular}{r@{\;}c}$\text{if }r=2:$ & \jepGraph{0/0/0/{2}/{}/0,1/1/0/{0}/{}/0,2/2.55/0/{0}/{}/0,3/3.55/0/{0}/{}/0,4/4.55/0/{0}/{}/0}{0/1,2/3}{1/2}{3/4/1}{}\end{tabular} & \jepGraph{0/0/0/{1}/{}/0,1/1/0/{0}/{}/0,2/2.55/0/{0}/{}/0,3/3.55/0/{0}/{}/0,4/4.55/0/{0}/{}/0}{0/1,2/3}{1/2}{3/4/1}{}\rule[-12pt]{0pt}{36pt}\\
\hline
\multirow{2}{*}{\begin{tabular}{c}$(C_n,\varpi_n)$\\[2pt]\jepGraph{0/0/0/{}/{1}/0,1/1.55/0/{}/{}/0,2/2.55/0/{}/{}/0,3/3.55/0/{}/{n}/1}{1/2}{0/1}{2/3/-1}{}\end{tabular}} & \multirow{2}{*}{\begin{tabular}{c}$\mathfrak{sp}(2n,\mathbb R)$\\[4pt]$r_M=n$\end{tabular}} & \begin{tabular}{r@{\;}c}$\text{if }r<n:$ & \jepGraph{0/0/0/{0}/{}/0,1/1.55/0/{0}/{}/0,2/2.55/0/{2}/{r}/0,3/3.55/0/{0}/{}/0,4/5.1/0/{0}/{}/0,5/6.1/0/{0}/{}/0,6/7.1/0/{0}/{}/0}{1/2,2/3,4/5}{0/1,3/4}{5/6/-1}{}\end{tabular} & \jepGraph{0/0/0/{0}/{}/0,1/1.55/0/{0}/{}/0,2/2.55/0/{2}/{r}/0,3/4.1/0/{2}/{}/0,4/5.1/0/{2}/{}/0,5/6.1/0/{1}/{}/0}{1/2,3/4}{0/1,2/3}{4/5/-1}{}\rule[-12pt]{0pt}{36pt}\\
\cline{3-4}
 & & \begin{tabular}{r@{\;}c}$\text{if }r=n:$ & \jepGraph{0/0/0/{0}/{}/0,1/1.55/0/{0}/{}/0,2/2.55/0/{0}/{}/0,3/3.55/0/{2}/{}/0}{1/2}{0/1}{2/3/-1}{}\end{tabular} & \jepGraph{0/0/0/{0}/{}/0,1/1.55/0/{0}/{}/0,2/2.55/0/{0}/{}/0,3/3.55/0/{1}/{}/0}{1/2}{0/1}{2/3/-1}{}\rule[-12pt]{0pt}{36pt}\\
\hline
\multirow{2}{*}{\begin{tabular}{c}$(D_n,\varpi_1)$\\[2pt]\jepGraph{0/0/0/{}/{1}/1,1/1/0/{}/{}/0,2/2.55/0/{}/{}/0,3/3.05/0.8/{}/{n-1}/0,4/3.05/-0.8/{}/{n}/0}{0/1,2/3,2/4}{1/2}{}{}\end{tabular}} & \multirow{2}{*}{\begin{tabular}{c}$\mathfrak{so}(2,2n-2)$\\[4pt]$r_M=2$\end{tabular}} & \begin{tabular}{r@{\;}c}$\text{if }r=1:$ & \jepGraph{0/0/0/{0}/{}/0,1/1/0/{1}/{}/0,2/2/0/{0}/{}/0,3/3.55/0/{0}/{}/0,4/4.55/0/{0}/{}/0,5/5.05/0.8/{0}/{}/0,6/5.05/-0.8/{0}/{}/0}{0/1,1/2,3/4,4/5,4/6}{2/3}{}{}\end{tabular} & \jepGraph{0/0/0/{1}/{}/0,1/1/0/{2}/{}/0,2/2.55/0/{2}/{}/0,3/3.55/0/{2}/{}/0,4/4.05/0.8/{1}/{}/0,5/4.05/-0.8/{1}/{}/0}{0/1,2/3,3/4,3/5}{1/2}{}{}\rule[-12pt]{0pt}{36pt}\\
\cline{3-4}
 & & \begin{tabular}{r@{\;}c}$\text{if }r=2:$ & \jepGraph{0/0/0/{2}/{}/0,1/1/0/{0}/{}/0,2/2.55/0/{0}/{}/0,3/3.55/0/{0}/{}/0,4/4.05/0.8/{0}/{}/0,5/4.05/-0.8/{0}/{}/0}{0/1,2/3,3/4,3/5}{1/2}{}{}\end{tabular} & \jepGraph{0/0/0/{1}/{}/0,1/1/0/{0}/{}/0,2/2.55/0/{0}/{}/0,3/3.55/0/{0}/{}/0,4/4.05/0.8/{0}/{}/0,5/4.05/-0.8/{0}/{}/0}{0/1,2/3,3/4,3/5}{1/2}{}{}\rule[-12pt]{0pt}{36pt}\\
\hline
\multirow{3}{*}{\begin{tabular}{c}$(D_n,\varpi_n)$\\[2pt]\jepGraph{0/0/0/{}/{1}/0,1/1/0/{}/{}/0,2/2.55/0/{}/{}/0,3/3.05/0.8/{}/{n-1}/0,4/3.05/-0.8/{}/{n}/1}{0/1,2/3,2/4}{1/2}{}{}\end{tabular}} & \multirow{3}{*}{\begin{tabular}{c}$\mathfrak{so}^{\star}(2n)$\\[4pt]$r_M=\lfloor n/2\rfloor$\end{tabular}} & \begin{tabular}{r@{\;}c}$\text{if }2r\leqslant n-2:$ & \jepGraph{0/0/0/{0}/{}/0,1/1.55/0/{0}/{}/0,2/2.55/0/{1}/{2r}/0,3/3.55/0/{0}/{}/0,4/5.1/0/{0}/{}/0,5/6.1/0/{0}/{}/0,6/6.6/0.8/{0}/{}/0,7/6.6/-0.8/{0}/{}/0}{1/2,2/3,4/5,5/6,5/7}{0/1,3/4}{}{}\end{tabular} & \jepGraph{0/0/0/{0}/{}/0,1/1.55/0/{0}/{}/0,2/2.55/0/{1}/{}/0,3/3.55/0/{2}/{2r}/0,4/5.1/0/{2}/{}/0,5/6.1/0/{2}/{}/0,6/6.6/0.8/{1}/{}/0,7/6.6/-0.8/{1}/{}/0}{1/2,2/3,4/5,5/6,5/7}{0/1,3/4}{}{}\rule[-12pt]{0pt}{36pt}\\
\cline{3-4}
 & & \begin{tabular}{r@{\;}c}$\text{if }2r=n-1:$ & \jepGraph{0/0/0/{0}/{}/0,1/1.55/0/{0}/{}/0,2/2.55/0/{0}/{}/0,3/3.05/0.8/{1}/{}/0,4/3.05/-0.8/{1}/{}/0}{1/2,2/3,2/4}{0/1}{}{}\end{tabular} & \jepGraph{0/0/0/{0}/{}/0,1/1.55/0/{0}/{}/0,2/2.55/0/{1}/{}/0,3/3.05/0.8/{1}/{}/0,4/3.05/-0.8/{1}/{}/0}{1/2,2/3,2/4}{0/1}{}{}\rule[-12pt]{0pt}{36pt}\\
\cline{3-4}
 & & \begin{tabular}{r@{\;}c}$\text{if }2r=n:$ & \jepGraph{0/0/0/{0}/{}/0,1/1.55/0/{0}/{}/0,2/2.55/0/{0}/{}/0,3/3.05/0.8/{0}/{}/0,4/3.05/-0.8/{2}/{}/0}{1/2,2/3,2/4}{0/1}{}{}\end{tabular} & \jepGraph{0/0/0/{0}/{}/0,1/1.55/0/{0}/{}/0,2/2.55/0/{0}/{}/0,3/3.05/0.8/{0}/{}/0,4/3.05/-0.8/{1}/{}/0}{1/2,2/3,2/4}{0/1}{}{}\rule[-12pt]{0pt}{36pt}\\
\hline
\multirow{2}{*}{\begin{tabular}{c}$(E_6,\varpi_1)$\\[2pt]\jepGraph{0/0/0/{}/{1}/1,1/1/0/{}/{}/0,2/2/0/{}/{}/0,3/3/0/{}/{}/0,4/4/0/{}/{}/0,5/2/1/{}/{}/0}{0/1,1/2,2/3,3/4,2/5}{}{}{}\end{tabular}} & \multirow{2}{*}{\begin{tabular}{c}$\mathfrak e_{6(-14)}$\\[4pt]$r_M=2$\end{tabular}} & \begin{tabular}{r@{\;}c}$\text{if }r=1:$ & \jepGraph{0/0/0/{0}/{}/0,1/1/0/{0}/{}/0,2/2/0/{0}/{}/0,3/3/0/{0}/{}/0,4/4/0/{0}/{}/0,5/2/1/{1}/{}/0}{0/1,1/2,2/3,3/4,2/5}{}{}{}\end{tabular} & \jepGraph{0/0/0/{1}/{}/0,1/1/0/{2}/{}/0,2/2/0/{3}/{}/0,3/3/0/{2}/{}/0,4/4/0/{1}/{}/0,5/2/1/{2}/{}/0}{0/1,1/2,2/3,3/4,2/5}{}{}{}\rule[-12pt]{0pt}{36pt}\\
\cline{3-4}
 & & \begin{tabular}{r@{\;}c}$\text{if }r=2:$ & \jepGraph{0/0/0/{1}/{}/0,1/1/0/{0}/{}/0,2/2/0/{0}/{}/0,3/3/0/{0}/{}/0,4/4/0/{1}/{}/0,5/2/1/{0}/{}/0}{0/1,1/2,2/3,3/4,2/5}{}{}{}\end{tabular} & \jepGraph{0/0/0/{1}/{}/0,1/1/0/{1}/{}/0,2/2/0/{1}/{}/0,3/3/0/{1}/{}/0,4/4/0/{1}/{}/0,5/2/1/{0}/{}/0}{0/1,1/2,2/3,3/4,2/5}{}{}{}\rule[-12pt]{0pt}{36pt}\\
\hline
\multirow{3}{*}{\begin{tabular}{c}$(E_7,\varpi_7)$\\[2pt]\jepGraph{0/0/0/{}/{}/0,1/1/0/{}/{}/0,2/2/0/{}/{}/0,3/3/0/{}/{}/0,4/4/0/{}/{}/0,5/5/0/{}/{7}/1,6/2/1/{}/{}/0}{0/1,1/2,2/3,3/4,4/5,2/6}{}{}{}\end{tabular}} & \multirow{3}{*}{\begin{tabular}{c}$\mathfrak e_{7(-25)}$\\[4pt]$r_M=3$\end{tabular}} & \begin{tabular}{r@{\;}c}$\text{if }r=1:$ & \jepGraph{0/0/0/{1}/{}/0,1/1/0/{0}/{}/0,2/2/0/{0}/{}/0,3/3/0/{0}/{}/0,4/4/0/{0}/{}/0,5/5/0/{0}/{}/0,6/2/1/{0}/{}/0}{0/1,1/2,2/3,3/4,4/5,2/6}{}{}{}\end{tabular} & \jepGraph{0/0/0/{2}/{}/0,1/1/0/{3}/{}/0,2/2/0/{4}/{}/0,3/3/0/{3}/{}/0,4/4/0/{2}/{}/0,5/5/0/{1}/{}/0,6/2/1/{2}/{}/0}{0/1,1/2,2/3,3/4,4/5,2/6}{}{}{}\rule[-12pt]{0pt}{36pt}\\
\cline{3-4}
 & & \begin{tabular}{r@{\;}c}$\text{if }r=2:$ & \jepGraph{0/0/0/{0}/{}/0,1/1/0/{0}/{}/0,2/2/0/{0}/{}/0,3/3/0/{0}/{}/0,4/4/0/{1}/{}/0,5/5/0/{0}/{}/0,6/2/1/{0}/{}/0}{0/1,1/2,2/3,3/4,4/5,2/6}{}{}{}\end{tabular} & \jepGraph{0/0/0/{0}/{}/0,1/1/0/{1}/{}/0,2/2/0/{2}/{}/0,3/3/0/{2}/{}/0,4/4/0/{2}/{}/0,5/5/0/{1}/{}/0,6/2/1/{1}/{}/0}{0/1,1/2,2/3,3/4,4/5,2/6}{}{}{}\rule[-12pt]{0pt}{36pt}\\
\cline{3-4}
 & & \begin{tabular}{r@{\;}c}$\text{if }r=3:$ & \jepGraph{0/0/0/{0}/{}/0,1/1/0/{0}/{}/0,2/2/0/{0}/{}/0,3/3/0/{0}/{}/0,4/4/0/{0}/{}/0,5/5/0/{2}/{}/0,6/2/1/{0}/{}/0}{0/1,1/2,2/3,3/4,4/5,2/6}{}{}{}\end{tabular} & \jepGraph{0/0/0/{0}/{}/0,1/1/0/{0}/{}/0,2/2/0/{0}/{}/0,3/3/0/{0}/{}/0,4/4/0/{0}/{}/0,5/5/0/{1}/{}/0,6/2/1/{0}/{}/0}{0/1,1/2,2/3,3/4,4/5,2/6}{}{}{}\rule[-12pt]{0pt}{36pt}\\
\hline
\end{tabular}
\caption{\small Dominant orthogonal sequences $(\alpha_1,\ldots,\alpha_r)$ for $1\leqslant r\leqslant r_M$}
\label{jep93table1}
\begin{minipage}{.92\textwidth}\footnotesize
(We do not indicate all the roots $\alpha_i$, but rather the sum of the corresponding coroots $\alpha_1^\jepPubCoroot+\cdots+\alpha_r^\jepPubCoroot$, by indicating on a Dynkin diagram the values $\langle\alpha_1^\jepPubCoroot+\cdots+\alpha_r^\jepPubCoroot,\beta\rangle$ for all simple roots $\beta$. In the last column, we express the root $\alpha_r$ which is also the smallest root $\alpha$ such that $\langle\alpha_1^\jepPubCoroot+\cdots+\alpha_r^\jepPubCoroot,\alpha\rangle=2$.)
\end{minipage}
\end{table}
\clearpage

\begin{proof}
Kostant \cite[Lem.~1.6]{jep93Kos12} proved the strong orthogonality. For the convenience of the reader, we recall his arguments in our context. The root $\alpha_i$ is the highest root of $R(\Pi_i)$, and $\Pi_i$ contains the long root $\zeta$, so $\alpha_i$ is long. By construction, $\langle\alpha_i^\jepPubCoroot,\alpha_j\rangle=0$ if $j>i$, so $\alpha_i$ and $\alpha_j$ are orthogonal. Since they are both long, they are strongly orthogonal, otherwise $\alpha_i\pm\alpha_j$ would be longer than $\alpha_i$. This proves~(1).

For~(2), by Fact~2.2, it is enough to prove that $(\alpha_1,\ldots,\alpha_q)$ is a maximal sequence of orthogonal roots (see also \cite[Th.~1.8]{jep93Kos12}). Let $\alpha\in R$ be such that $\langle\alpha_i^\jepPubCoroot,\alpha\rangle=0$ holds for all $i$. We can and will assume that $\alpha>0$. Let $i$ be the greatest integer such that $\alpha\in R(\Pi_i)$. By maximality of $i$ there exists a simple root $\beta$ in $\operatorname{supp}(\alpha)\cap\Sigma_i$. Since $\alpha_i$ is dominant on $R(\Pi_i)$, we have $\langle\alpha_i^\jepPubCoroot,\alpha\rangle\geqslant\langle\alpha_i^\jepPubCoroot,\beta\rangle>0$, a contradiction to the assumption made on $\alpha$.

For the third point, let $i\leqslant r_M$ and let $\beta\in\Pi$.
\begin{itemize}
\item If $\beta\in\Pi_i$, by construction, $\langle\alpha_j^\jepPubCoroot,\beta\rangle=0$ for $j<i$. Since $\alpha_i$ is the highest root of $\Pi_i$, $\langle\alpha_i^\jepPubCoroot,\beta\rangle\geqslant0$. Thus $\langle\alpha_1^\jepPubCoroot+\cdots+\alpha_i^\jepPubCoroot,\beta\rangle=\langle\alpha_i^\jepPubCoroot,\beta\rangle\geqslant0$
\item If $\beta\in\Sigma_{i-1}$, then $\langle\alpha_{i-1}^\jepPubCoroot,\beta\rangle\geqslant1$ and $\langle\alpha_j^\jepPubCoroot,\beta\rangle=0$ for $j<i-1$. Since $\alpha_i$ is long, $\langle\alpha_i^\jepPubCoroot,\beta\rangle\geqslant-1$. Thus, $\langle\alpha_1^\jepPubCoroot+\cdots+\alpha_i^\jepPubCoroot,\beta\rangle\geqslant0$
\item If $\beta\in\Sigma_j$ for $j<i-1$, then, by construction of $\Pi_i$, $\langle\alpha_i^\jepPubCoroot,\beta\rangle=0$. By induction on $i$, $\langle\alpha_1^\jepPubCoroot+\cdots+\alpha_{i-1}^\jepPubCoroot,\beta\rangle\geqslant0$, so $\langle\alpha_1^\jepPubCoroot+\cdots+\alpha_i^\jepPubCoroot,\beta\rangle\geqslant0$
\end{itemize}
For~(4), by Lemma~2.8 and since $\alpha_i$ is long, we have $\langle\varpi,\alpha_i^\jepPubCoroot\rangle=1$. Thus $s_{\alpha_i}(\varpi)=\varpi-\langle\varpi,\alpha_i^\jepPubCoroot\rangle\alpha_i=\varpi-\alpha_i$. Therefore $s_{\alpha_1}\cdots s_{\alpha_{r_M}}(\varpi)=\varpi-\alpha_1-\cdots-\alpha_{r_M}$ is a weight of $\jepPubBbb{E}$. We prove that it is a lowest weight. Let $\beta\in\Pi$. If $\beta\neq\zeta$, then
\[
\langle\varpi-\alpha_1-\cdots-\alpha_{r_M},\beta^\jepPubCoroot\rangle=\langle-\alpha_1-\cdots-\alpha_{r_M},\beta^\jepPubCoroot\rangle,
\]
and this is non-positive by~(3) and because all the roots $\alpha_i$ have the same length. If $\beta=\zeta$, then we compute that $\langle\varpi-\alpha_1-\cdots-\alpha_{r_M},\zeta^\jepPubCoroot\rangle=1-\langle\alpha_{r_M},\zeta^\jepPubCoroot\rangle\leqslant0$ since $\zeta\in\Sigma_{r_M}$.
\end{proof}
We make the following observations.
\begin{proposition}
Let $1\leqslant r\leqslant r_M$, $(\alpha_1,\ldots,\alpha_r)$ be the corresponding dominant orthogonal sequence, and $h_r=\alpha_1^\jepPubCoroot+\cdots+\alpha_r^\jepPubCoroot$. Then:
\begin{enumerate}[label=(\arabic*)]
\item $h_r$ is dominant: for any positive root $\alpha$, we have $\langle\alpha,h_r\rangle\geqslant0$
\item For any root $\alpha$, we have $\langle\alpha,h_r\rangle\leqslant2$
\item $\alpha_r$ is the smallest root $\alpha$ such that $\langle\alpha,h_r\rangle=2$
\item If a root $\alpha$ satisfies $\langle\alpha,h_r\rangle=2$, then $\langle\alpha,z\rangle=2$
\item We have $\langle\varpi,h_r\rangle=r$
\end{enumerate}
\end{proposition}
\begin{proof}
Recall Table~1. The first point has been proved in Proposition~2.19~(3). Let $\Theta$ be the highest root of the root system of $G$, which can be found for example in \cite{jep93Bou68}. Since $h_r$ is dominant, the second item follows from the fact that $\langle\Theta,h_r\rangle=2$ in all cases.

For the third item, we have indicated the root $\alpha_r$ in the last column of the table. It readily follows that it is the smallest root $\alpha$ such that $\langle\alpha,h_r\rangle=2$. Since it has coefficient $1$ on $\zeta$, we have $\langle\alpha_r,z\rangle=2$. Any root $\alpha$ as in the fourth item will be bigger than $\alpha_r$, so the claim follows.

To prove that $\langle\varpi,h_r\rangle=r$, recall that if $(\alpha_1,\ldots,\alpha_{r_M})$ is the maximal dominant orthogonal sequence, the weight $\varpi-\sum_{1}^{r_M}\alpha_i$ is the lowest weight $\mu_0$ by Proposition~2.19~(4). Thus, $\langle\varpi-\mu_0,h_r\rangle=\langle\alpha_1+\cdots+\alpha_{r_M},\alpha_1^\jepPubCoroot+\cdots+\alpha_r^\jepPubCoroot\rangle=2r$, since $\langle\alpha_i,\alpha_j^\jepPubCoroot\rangle=2\delta_{i,j}$ by orthogonality of the roots $\alpha_i$. Moreover, $h_r$ is part of some $\mathfrak{sl}_2$-triple $(x,h_r,y)$, so $\langle\mu_0,h_r\rangle=-\langle\varpi,h_r\rangle$, by the representation theory of $\mathfrak{sl}_2$. This proves that $\langle\varpi,h_r\rangle=r$.
\end{proof}

